%
% integral.tex -- Integration der geometrischen Reihe
%
% (c) 2021 Prof Dr Andreas Müller, OST Ostschweizer Fachhochschule
%
\documentclass[tikz]{standalone}
\usepackage{amsmath}
\usepackage{times}
\usepackage{txfonts}
\usepackage{pgfplots}
\usepackage{csvsimple}
\definecolor{darkred}{rgb}{0.8,0,0}
\usetikzlibrary{arrows,intersections,math,calc}
\begin{document}
\def\skala{1}
\pgfmathparse{atan(4/3)}
\xdef\w{\pgfmathresult}
\def\punkt#1#2{
	\fill[color=white] #1 circle[radius=0.05];
	\draw[color=#2] #1 circle[radius=0.05];
}
\begin{tikzpicture}[>=latex,thick,scale=\skala]

\coordinate (A) at (3,3);
\coordinate (c) at (0.9,1.2);
\coordinate (d) at (-1,-0.5);

\fill[color=gray!10] (A) circle[radius=1.5];
\draw[->,line width=0.6pt] (A) -- ($(A)+(c)$);
\node at ($(A)+0.7*(c)$) [left] {$\varrho$};
\draw[color=blue,line width=0.4pt] (A) -- ++(d);

\node at (A) [below] {$z_0$};
\node[color=blue] at ($(A)+(d)$) [below] {$z$};

\coordinate (B1) at (2,0.8);
\coordinate (B2) at (5,2);
\coordinate (B3) at (1,4);

\draw[color=darkred,line width=1.2pt] (B1) to[out=135,in=-120]
	(B2) to[out=60,in={\w-90}]
	($(A)+(c)$) to[out={90+\w},in=90]
	(B3);

\punkt{(A)}{black}
\punkt{($(A)+(c)$)}{darkred}
\punkt{($(A)+(d)$)}{blue}

\node[color=darkred] at (4.7,3.6) [right] {$x=\gamma(t)$};

\draw[->] (-0.1,0) -- (6.3,0) coordinate[label={$\operatorname{Re}$}];
\draw[->] (0,-0.1) -- (0,5.3) coordinate[label={right:$\operatorname{Im}$}];

\end{tikzpicture}
\end{document}

